\newcommand{\TabOptimisationParameters}{%
\begin{table}[tb]%
\begin{tabular}{ll}%
	\hline
	Region for optimisation & Approx.\ No.\  of parameters \\
	\hline
	Production target dimensions and location & $3+3$ \\
	Torus1 dipole field strength & 1 \\
	Torus2 dipole field strength & 1 \\
	Muon beam collimator shapes, position, and material & $3+1+1$ \\
	Stopping target shape and location & $4+3$ \\
	Beam blocker position, form, and material& $3+3+1$ \\
	Electron spectrometer dipole field strength& 1 \\
	DIO blockers in the spectrometer & $4$ \\
	\hline
	Approx. total number of parameters& 32 \\
	\hline
\end{tabular}%
	\caption{\tablabel{optimisation:possible-parameters}%
	Aspects of the experiment that can be optimised and estimates for the number of parameters that define each aspect.
	In the case of the target, beam blocker, and collimator shapes the number of parameters is only approximate; crudely speaking there is at least a width, length and height but in principle one could have a very irregular shape that cannot be parametrised by only three numbers, for example, shapes that change as a function of distance along the beamline.
}%
\end{table}%
}

\newcommand{\FigOptimProdTgtLength}{
\begin{figure}[pt]
\centering
\subfloat[][\figlabel{optimisation:ProdTgtSec:Length:Momentum:Muons}Muons]{
\includegraphics[width=0.48\textwidth,trim=0 0 0 1.5cm,clip]{figs/optimisation/ProdTgtGeom/Length_mu-minus_momentum}}
\subfloat[][\figlabel{optimisation:ProdTgtSec:Length:Momentum:Pions}Pions]{
\includegraphics[width=0.48\textwidth,trim=0 0 0 1.5cm,clip]{figs/optimisation/ProdTgtGeom/Length_pi-minus_momentum}}
\caption{\figlabel{optimisation:ProdTgtSec:Length:Momentum}
Change to momentum distributions at the entrance to the first 90 degrees of the bent muon beam solenoid for different target lengths.
Target length is given as half-length which is the Geant4 convention.  
}
\end{figure}
\begin{figure}[pt]
\centering
\subfloat[][\figlabel{optimisation:ProdTgtSec:Length:Integral:Muons}Muons]{
\includegraphics[width=0.48\textwidth,trim=0 0 2cm 2.8cm,clip]{figs/optimisation/ProdTgtGeom/Length_mu-minus_integral_toZero}}
\subfloat[][\figlabel{optimisation:ProdTgtSec:Length:Integral:Pions}Pions]{
\includegraphics[width=0.48\textwidth,trim=0 0 2cm 2.8cm,clip]{figs/optimisation/ProdTgtGeom/Length_pi-minus_integral_toZero}}
\caption{\figlabel{optimisation:ProdTgtSec:Length:Integral}
Integrated muon and pion yields up to a certain momentum at the entrance to the first 90 degrees of the bent muon beam solenoid as a function of target length.
}
\end{figure}
\begin{figure}[pt]
\centering
\subfloat[][\figlabel{optimisation:ProdTgtSec:Length:IntegralRatio:Muons}Muons]{\includegraphics[width=0.48\textwidth,trim=0 0 2cm 2.8cm,clip]{figs/optimisation/ProdTgtGeom/Length_mu-minus_integral_ratios}}
\subfloat[][\figlabel{optimisation:ProdTgtSec:Length:IntegralRatio:Pions}Pions]{\includegraphics[width=0.48\textwidth,trim=0 0 2cm 2.8cm,clip]{figs/optimisation/ProdTgtGeom/Length_pi-minus_integral_ratios}}
\caption{\figlabel{optimisation:ProdTgtSec:Length:IntegralRatio}
Change in the momentum distribution of muons and pions at the entrance to the first 90 degrees of the bent muon beam solenoid as a function of target length.
}
\end{figure}
}

\newcommand{\FigOptimProdTgtRad}{
\begin{figure}[pt]
\centering
\subfloat[][\figlabel{optimisation:ProdTgtSec:Radius:Momentum:Muons}Muons]{\includegraphics[width=0.48\textwidth,trim=0 0 0 1.5cm,clip]{figs/optimisation/ProdTgtGeom/Radius_mu-minus_momentum}}
\subfloat[][\figlabel{optimisation:ProdTgtSec:Radius:Momentum:Pions}Pions]{\includegraphics[width=0.48\textwidth,trim=0 0 0 1.5cm,clip]{figs/optimisation/ProdTgtGeom/Radius_pi-minus_momentum}}
\caption{
\figlabel{optimisation:ProdTgtSec:Radius:Momentum}
Change to momentum distributions at the entrance to the first 90 degrees of the bent muon beam solenoid for different target radii.
}
\end{figure}
\begin{figure}[pt]
\centering
\subfloat[][\figlabel{optimisation:ProdTgtSec:Radius:Integral:Muons}Muons]{\includegraphics[width=0.48\textwidth,trim=0 0 2cm 2.8cm,clip]{figs/optimisation/ProdTgtGeom/Radius_mu-minus_integral_toZero}}
\subfloat[][\figlabel{optimisation:ProdTgtSec:Radius:Integral:Pions}Pions]{\includegraphics[width=0.48\textwidth,trim=0 0 2cm 2.8cm,clip]{figs/optimisation/ProdTgtGeom/Radius_pi-minus_integral_toZero}}
\caption{\figlabel{optimisation:ProdTgtSec:Radius:Integral}
Integrated muon and pion yields up to a certain momentum at the entrance to the first 90 degrees of the bent muon beam solenoid as a function of target radius.
}
\end{figure}
\begin{figure}[pt]
\centering
\subfloat[][\figlabel{optimisation:ProdTgtSec:Radius:IntegralRatio:Muons}Muons]{\includegraphics[width=0.48\textwidth,trim=0 0 2cm 2.8cm,clip]{figs/optimisation/ProdTgtGeom/Radius_mu-minus_integral_ratios}}
\subfloat[][\figlabel{optimisation:ProdTgtSec:Radius:IntegralRatio:Pions}Pions]{\includegraphics[width=0.48\textwidth,trim=0 0 2cm 2.8cm,clip]{figs/optimisation/ProdTgtGeom/Radius_pi-minus_integral_ratios}}
\caption{\figlabel{optimisation:ProdTgtSec:Radius:IntegralRatio}
Change in the momentum distribution of muons and pions at the entrance to the first 90 degrees of the bent muon beam solenoid as a function of target radius.
}
\end{figure}
}

\newcommand{\FigOptimProdTgtFinal}{
\begin{figure}[pt]
\centering
\subfloat[][\figlabel{optimisation:ProdTgtSec:Final:Integral:Muons}Muons]{\includegraphics[width=0.48\textwidth,trim=0 0 0 1.5cm,clip]{figs/optimisation/ProdTgtGeom/OptimalLengthRadius_mu-minus_integral_toZero.pdf}}
\subfloat[][\figlabel{optimisation:ProdTgtSec:Final:Integral:Pions}Pions]{\includegraphics[width=0.48\textwidth,trim=0 0 0 1.5cm,clip]{figs/optimisation/ProdTgtGeom/OptimalLengthRadius_pi-minus_integral_toZero.pdf}}
\caption{\figlabel{optimisation:ProdTgtSec:Final:Integral}
Variation in muon and pion yields as a function of target radius when the total target length is set to the optimised value of 32~cm.
Despite the longer target length the optimal radius is still 1~cm.
}
\end{figure}
}

\newcommand{\FigOptimProdTgtComparePhases}{
\begin{figure}[pt]
\centering
\includegraphics[width=0.9\textwidth,trim=0 0 0 1.5cm,clip]{figs/optimisation/ProdTgtGeom/Plot_compare_phase_1and2.pdf}
\caption{\figlabel{optimisation:ProdTgtSec:Phase1vs2}
Comparison of the muon and pion yields per POT for \phaseI and \phaseII.
The difference arises from the change of target material between the phases.
}
\end{figure}
}
\newcommand{\FigOptimMuBeamDipoleMuStops}{
\begin{figure}[bt]
\centering
%	\fbox{
\includegraphics[width=0.85\textwidth,trim=0 0.5cm 0 1.0cm,clip]{figs/optimisation/MuonBeamDipoles/Tidied_stopped_muons.pdf}
%}
\caption{\figlabel{optim:muBeamDipole:stoppedMu}
	Muon stopping rate as a function of the two dipole field strengths (given relative to the \phaseI design specification).
	A clear anti-correlation is visible which is discussed in the text.
}
\end{figure}
}

\newcommand{\FigOptimMuBeamDipolePiStops}{
\begin{figure}[t]
\centering
%	\fbox{
\includegraphics[width=0.85\textwidth,trim=0 0.5cm 0 1.0cm,clip]{figs/optimisation/MuonBeamDipoles/Tidied_stopped_pions.pdf}
%}
\caption{\figlabel{optim:muBeamDipole:stoppedPi}
	Pion stopping rate as a function of the two dipole field strengths (given relative to the \phaseI design specification).
At the level of statistics used to generate each point, no clear trend is obvious.
Empty squares are those where no pions stopped in the run.
}
\end{figure}
}

\newcommand{\FigOptimMuBeamDipoleMuDispersion}{
\begin{figure}[t]
\centering 
\subfloat[][\figlabel{optim:MuBeamDipole:MuDispersion:Entry}Torus1 Entry]{\includegraphics[width=0.45\textwidth,trim=0 0.9cm 0 1.9cm,clip]{figs/optimisation/MuonBeamCollimators/Tidied_Muon_dispersion_at_entrance.pdf}}
\subfloat[][\figlabel{optim:MuBeamDipole:MuDispersion:Exit}Torus2 Exit]  {\includegraphics[width=0.45\textwidth,trim=0 0.9cm 0 1.9cm,clip]{figs/optimisation/MuonBeamCollimators/Tidied_Muon_dispersion_at_exit.pdf}}
\caption{\figlabel{optim:MuBeamDipole:MuDispersion}
Dispersive effect of the 180\degree bent transport solenoid and dipole field on muons.
No collimating material is yet included, so the high-energy muons being removed is due purely to the beam-pipe itself.
}
\end{figure}
}

\newcommand{\FigOptimMuBeamCollimMuonPaths}{
\begin{figure}[p]
\centering 
	\subfloat[][\figlabel{optim:MuBeamCollim:Beamline:All}All Muons]          {\includegraphics[width=1\textwidth,trim=18cm 1.0cm 26cm 1cm,clip]{figs/optimisation/MuonBeamCollimators/Tidied_All_muons_wGeom.png}}\\
\subfloat[][\figlabel{optim:MuBeamCollim:Beamline:Stopped}Stopped Muons]          {\includegraphics[width=1\textwidth,trim=18cm 1.0cm 26cm 1cm,clip]{figs/optimisation/MuonBeamCollimators/Tidied_Stopped_muons_wGeom.png}}\\
\subfloat[][\figlabel{optim:MuBeamCollim:Beamline:HighP}Muons with $p>70$ MeV/c around the stopping target]          {\includegraphics[width=1\textwidth,trim=18cm 1.0cm 26cm 1cm,clip]{figs/optimisation/MuonBeamCollimators/Tidied_HighP_muons_wGeom.png}}\\
\subfloat[][\figlabel{optim:MuBeamCollim:Beamline:Diff}Stopped $\mu-$High-$p$ $\mu$]{\includegraphics[width=1\textwidth,trim=18cm 1.0cm 26cm 1cm,clip]{figs/optimisation/MuonBeamCollimators/Tidied_Where_to_collimate_wGeom.png}}
\caption{\figlabel{optim:MuBeamCollim:Beamline}
The heights of muons as they pass along the beamline.  
	\protect\subref{fig:optim:MuBeamCollim:Beamline:All} The path of all muons.
	\protect\subref{fig:optim:MuBeamCollim:Beamline:Stopped}: The paths of muons that stop in the target.
	\protect\subref{fig:optim:MuBeamCollim:Beamline:HighP}: The heights of muons with momentum greater than 70 MeV/c when they enter the region around the stopping target.  These could potentially decay in flight to give electrons with 100 MeV/c or greater.
	\protect\subref{fig:optim:MuBeamCollim:Beamline:Diff}: The difference between plot \protect\subref{fig:optim:MuBeamCollim:Beamline:Stopped} and plot \protect\subref{fig:optim:MuBeamCollim:Beamline:HighP}.
	Regions in dark blue would give the greatest impact in removing high-momentum muons whilst leave the stopping muons untouched.
	These plots should be compared to those of \fig{optim:MuBeamCollim:BeamWColl} once collimators have been introduced.
}
\end{figure}
}

\newcommand{\FigOptimMuBeamCollimTransverseSep}{
\begin{figure}[t]
\centering 
\subfloat[][\figlabel{optim:MuBeamCollim:TransverseSep:TS1Entry}At 0\degree]  {\includegraphics[height=0.35\textheight,trim=0.0cm 0.8cm 1.3cm 1.9cm,clip]{figs/optimisation/MuonBeamCollimators/MuonTransversePos_Torus1Tor1FirstPoint}}
\subfloat[][\figlabel{optim:MuBeamCollim:TransverseSep:TS3}At 90\degree] {\includegraphics[height=0.35\textheight,trim=1.7cm 0.8cm 1.3cm 1.9cm,clip]{figs/optimisation/MuonBeamCollimators/MuonTransversePos_Torus2Monitor_0}}
\subfloat[][\figlabel{optim:MuBeamCollim:TransverseSep:TS2Exit}At 180\degree]{\includegraphics[height=0.35\textheight,trim=1.7cm 0.8cm 1.3cm 1.9cm,clip]{figs/optimisation/MuonBeamCollimators/MuonTransversePos_StopTgtSecMonitor_0}}
\caption{\figlabel{optim:MuBeamCollim:TransverseSep}
The separation between stopping and dangerous muons.
The separation is largest at the exit (180\degree), reasonable at the entrance (0\degree), and smallest around the mid-point (90\degree).
}
\end{figure}
}

\newcommand{\FigOptimMuBeamCollimMuonPathsWColl}{
\begin{figure}[ph]
\centering 
\subfloat[][\figlabel{optim:MuBeamCollim:BeamWColl:All}All Muons]                  {\includegraphics[width=1\textwidth,trim=18cm 1.0cm 26cm 1cm,clip]{figs/optimisation/MuonBeamCollimators/Tidied_WColl_AllMuons_WGeom.png}}\\
\subfloat[][\figlabel{optim:MuBeamCollim:BeamWColl:Stopped}Stopped Muons]          {\includegraphics[width=1\textwidth,trim=18cm 1.0cm 26cm 1cm,clip]{figs/optimisation/MuonBeamCollimators/Tidied_WColl_StoppedMuons_WGeom.png}}\\
\subfloat[][\figlabel{optim:MuBeamCollim:BeamWColl:HighP}Muons with $p>70$ MeV/c around the stopping target]          {\includegraphics[width=1\textwidth,trim=18cm 1.0cm 26cm 1cm,clip]{figs/optimisation/MuonBeamCollimators/Tidied_WColl_HighPMuons_WGeom.png}}\\
\caption{\figlabel{optim:MuBeamCollim:BeamWColl}
The heights of muons as they pass along the beamline.  
	\protect\subref{fig:optim:MuBeamCollim:Beamline:All} The path of all muons.
	\protect\subref{fig:optim:MuBeamCollim:Beamline:Stopped}: The paths of muons that stop in the target.
	\protect\subref{fig:optim:MuBeamCollim:Beamline:HighP}: The heights of muons with momentum greater than 70 MeV/c when they enter the region around the stopping target.  These could potentially decay in flight to give electrons with 100 MeV/c or greater.
	These plots should be compared to those of \fig{optim:MuBeamCollim:Beamline} before collimators were introduced, where it is clear how well the dangerous muons are being suppressed.
}
\end{figure}
}

\newcommand{\FigOptimMuBeamCollimTorusOne}{
\begin{figure}[t]
\centering 
\subfloat[][\figlabel{optim:MuBeamCollim:Torus1:perPOT}Particles Surviving per POT]{\includegraphics[width=0.5\textwidth,trim=0.8cm 0.8cm 0.6cm 1.9cm,clip]{figs/optimisation/MuonBeamCollimators/Survived_Coll1_unNormalised-log.pdf}}
\subfloat[][\figlabel{optim:MuBeamCollim:Torus1:fraction}Fraction Surviving Collimator]{\includegraphics[width=0.5\textwidth,trim=0.8cm 0.8cm 0.6cm 1.9cm,clip]{figs/optimisation/MuonBeamCollimators/Survived_Coll1_Normalised-lin.pdf}}
\caption{\figlabel{optim:MuBeamCollim:Torus1}
The effect of changing the height of the collimator in Torus1 on the particle distributions.
}
\end{figure}
}

\newcommand{\FigOptimMuBeamCollimTorusTwoPerPOT}{
\begin{figure}[t]
\centering 
\includegraphics[width=0.95\textwidth,trim=0.3cm 0.3cm 0.8cm 0.1cm,clip]{figs/optimisation/MuonBeamCollimators/Survived_Coll2_unNormalised-lin.pdf}
\caption{\figlabel{optim:MuBeamCollim:Torus2:perPOT}
The number of particles reaching the end of the Torus2 solenoid per POT for different heights of both collimators in Torus1 and Torus2.
}
\end{figure}
}


\newcommand{\FigOptimMuBeamCollimTorusTwoFraction}{
\begin{figure}[t]
\centering 
\includegraphics[width=0.8\textwidth,trim=0.3cm 0.3cm 0.8cm 0.1cm,clip]{figs/optimisation/MuonBeamCollimators/Survived_Coll2_unNormalised-lin.pdf}
\caption{\figlabel{optim:MuBeamCollim:Torus2:fraction}
	The number of particles reaching the end of the Torus2 solenoid relative to the number that enter the Torus1 solenoid (\ie the survival probability) for different heights of both collimators in Torus1 and Torus2.
}
\end{figure}
}

\newcommand{\FigOptimMuBeamCollimTorusTwoContours}{
\begin{figure}[bt]
\centering 
%\fbox{
\includegraphics[width=0.75\textwidth,trim=0.1cm 0.3cm 2.8cm 1.1cm,clip]{figs/optimisation/MuonBeamCollimators/Survived_Coll2_StoppedVsHighP-Muons.pdf}
%}
\caption{\figlabel{optim:MuBeamCollim:Torus2:contours}
Contours showing 2.5 percentage point changes to the stopping (blue) and dangerous (red) muon flux, as a function of the collimator heights.
100\% acceptance is found in the bottom right corner.
%For example, for collimator heights within the first blue contour towards the bottom-right corner, less than 2.5\% of stopped muons are lost.
}
\end{figure}
}

\newcommand{\FigOptimESTDipoleBeamHeightTwoD}{
\begin{figure}[tp]
\centering 
	\subfloat[][\figlabel{optim:ESTDipole:Beam:0.0}No Dipole]  {\includegraphics[width=1\textwidth,trim=4cm 0.5cm 9.5cm 0.7cm,clip]{figs/optimisation/EST_dipole/Tidied_signal_height-dipole_00}}\\
\subfloat[][\figlabel{optim:ESTDipole:Beam:0.1}0.1~T Dipole]{\includegraphics[width=1\textwidth,trim=4cm 0.5cm 9.5cm 0.7cm,clip]{figs/optimisation/EST_dipole/Tidied_signal_height-dipole_10}}\\
\subfloat[][\figlabel{optim:ESTDipole:Beam:0.2}0.2~T Dipole]{\includegraphics[width=1\textwidth,trim=4cm 0.5cm 9.5cm 0.7cm,clip]{figs/optimisation/EST_dipole/Tidied_signal_height-dipole_20}}\\
\caption{\figlabel{optim:ESTDipole:Beam}
The heights of signal electrons for different dipole field values.
}
\end{figure}
}

\newcommand{\FigOptimESTDipoleBeamHeightMean}{
\begin{figure}[tb]
\centering 
%	\fbox{
\includegraphics[width=0.8\textwidth,trim=1.0cm 1.3cm 3.8cm 0.9cm,clip]{figs/optimisation/EST_dipole/Tidied_NoShift-Height}
%}
\caption{\figlabel{optim:ESTDipole:MeanHeight}
Mean height of signal electrons for different values of the dipole field strength.
}
\end{figure}
}

\newcommand{\FigOptimESTDipoleBeamFluxMean}{
\begin{figure}[tb]
\centering 
\includegraphics[width=0.8\textwidth,trim=1.0cm 0.3cm 3.8cm 1.8cm,clip]{figs/optimisation/EST_dipole/Tidied_NoShift-Flux}
\caption{\figlabel{optim:ESTDipole:MeanFlux}
Survival probability for signal electrons as a function of the distance along the beamline for different values of the electron spectrometer's dipole field strengths.
}
\end{figure}
}

\newcommand{\FigOptimESTDipoleAcceptanceVsDipole}{
\begin{figure}[tb]
\centering 
\includegraphics[width=0.6\textwidth,trim=0.3cm 0.3cm 1.9cm 1.1cm,clip]{figs/optimisation/EST_dipole/Tidied_acceptance}
\caption{\figlabel{optim:ESTDipole:acceptance}
Geometric acceptance into the StrECAL detector as a function of the dipole field strength over the electron spectrometer.
}
\end{figure}
}

\newcommand{\FigOptimStopTgtPosMuStops}{
\begin{figure}[bt]
\centering 
%	\fbox{
\includegraphics[width=0.7\textwidth,trim=2.0cm 0.05cm 0.9cm 0.3cm,clip]{figs/optimisation/StopTgtPosition/Tidied_MuonStoppingRate.pdf}
%}
\caption{\figlabel{optim:StopTgtPos:MuStops}
Muon stopping rate per \ac{POT} for different target positions.
The linear behaviour arises from the reduced field strength and fixed target radius such that fewer muons impact the target as it is moved downstream.
}
\end{figure}
}

\newcommand{\FigOptimStopTgtPosSensitivitySpect}{
\begin{figure}[tb]
\centering 
%	\fbox{
\includegraphics[width=0.9\textwidth,,trim=1.3cm 0.2cm 0.7cm 0.6cm,clip]{figs/optimisation/StopTgtPosition/Tidied_-AcceptedMomentum.pdf}
%}
\caption{\figlabel{optim:StopTgtPos:AcceptedMomSpect}
The momentum dependence of the electron acceptance into the detector for different target positions.
The spectrum for each target position is normalised to the muon stopping rate for that position, such that each curve shows the sensitivity to electrons of that momentum.
}
\end{figure}
}

\newcommand{\FigOptimStopTgtPosSensitivityNoBeamBlock}{
\begin{figure}[tb]
\centering 
%	\fbox{
\includegraphics[width=0.9\textwidth,trim=1.3cm 0.2cm 0.7cm 0.6cm,clip]{figs/optimisation/StopTgtPosition/Tidied_NoBeamBlocker-AcceptedMomentum.pdf}
%}
\caption{\figlabel{optim:StopTgtPos:AcceptedMomSpectNoBeamBlock}
The momentum dependence of the electron acceptance into the detector for different target positions when the beam blocker is removed.
}
\end{figure}
}

\newcommand{\FigOptimStopTgtPosSensitivityIntegral}{
\begin{figure}[tb]
\centering 
%	\fbox{
\subfloat[][\figlabel{optim:StopTgtPos:AcceptIntegral:Adjacent}Sensitivity]{%
\includegraphics[width=0.49\textwidth,trim=0.5cm 0.5cm 0.3cm 1.9cm,clip]{figs/optimisation/StopTgtPosition/Tidied_-AcceptedMomentum-Integrated_adjacent.pdf}}
\subfloat[][\figlabel{optim:StopTgtPos:AcceptIntegral:Ratio}Change in Shape]{%
\includegraphics[width=0.49\textwidth,trim=0.5cm 0.5cm 0.3cm 1.9cm,clip]{figs/optimisation/StopTgtPosition/Tidied_-AcceptedMomentum-Integrated_ratios.pdf}}
%}
\caption{\figlabel{optim:StopTgtPos:AcceptIntegral}
\protect\subref{fig:optim:StopTgtPos:AcceptIntegral:Adjacent} 
The variation in sensitivity (acceptance $\times$ stopping rate) to electrons with different momenta as a function of the target position with respect to the nominal location. 
The darkest red line towards the top of the plot represents the sensitivity to signal, and it is that line that should therefore be maximised.
\protect\subref{fig:optim:StopTgtPos:AcceptIntegral:Ratio} 
The change in the shape of the acceptance vs. momentum spectrum as a function of the stopping target location.
}
\end{figure}
}

\newcommand{\FigOptimStopTgtPosHeights}{
\begin{figure}[p]
\centering 
\subfloat[][\figlabel{optim:StopTgtPos:Height:42.5}Electrons from 40 to 45 MeV/c]{%
\includegraphics[width=0.95\textwidth,trim=1.15cm 0.05cm 0.3cm 0.180cm,clip]{figs/optimisation/StopTgtPosition/WithBeamBlocker-Height-VaryShifts-Momentum_42-5.pdf}}%
\\\subfloat[][\figlabel{optim:StopTgtPos:Height:62.5}Electrons from 60 to 65 MeV/c]{%
\includegraphics[width=0.95\textwidth,trim=1.15cm 0.05cm 0.3cm 0.180cm,clip]{figs/optimisation/StopTgtPosition/WithBeamBlocker-Height-VaryShifts-Momentum_62-5.pdf}}%
\\\subfloat[][\figlabel{optim:StopTgtPos:Height:82.5}Electrons from 80 to 85 MeV/c]{%
\includegraphics[width=0.95\textwidth,trim=1.15cm 0.05cm 0.3cm 0.180cm,clip]{figs/optimisation/StopTgtPosition/WithBeamBlocker-Height-VaryShifts-Momentum_82-5.pdf}}%
\\\subfloat[][\figlabel{optim:StopTgtPos:Height:102.5}Electrons from 100 to 105 MeV/c]{%
\includegraphics[width=0.95\textwidth,trim=1.15cm 0.05cm 0.3cm 0.180cm,clip]{figs/optimisation/StopTgtPosition/WithBeamBlocker-Height-VaryShifts-Momentum_102-5.pdf}}%
\\\subfloat[][\figlabel{optim:StopTgtPos:Height:102.5-zoom}Electrons from 100 to 105 MeV/c (Zoomed)]{%
\includegraphics[width=0.95\textwidth,trim=1.15cm 0.05cm 0.3cm 0.180cm,clip]{figs/optimisation/StopTgtPosition/WithBeamBlocker-Height-VaryShifts-Momentum_102-5-Zoom.pdf}}%
\caption{\figlabel{optim:StopTgtPos:Height}
The effect of stopping target position on the height of electrons with a fixed momentum as they pass through the Electron Spectrometer.
The size of the variation indicates the stability of the dipole tune; the two parameters are clearly correlated.
Also striking---particularly in \protect\subref{fig:optim:StopTgtPos:Height:102.5-zoom}---is the way the dependence on the helical pitch angles is affected by the stopping target position.
}
\end{figure}
}

%\newcommand{\FigOptimStopTgtPosSensitivityIntegral}{
%\begin{figure}[tb]
%\centering 
%%	\fbox{
%\includegraphics[width=0.9\textwidth,trim=1.2cm 0.05cm 0.9cm 0.6cm,clip]{figs/optimisation/StopTgtPosition/Tidied_-AcceptedMomentum-Integrated_adjacent.pdf}
%%}
%\caption{\figlabel{optim:StopTgtPos:AcceptedIntegrated}
%The variation in sensitivity to different momentum electrons as a function of the target position, with respect to the nominal location.
%The darkest red line towards the top of the plot represents the signal sensitivity, and it is that line that should therefore be maximised.
%}
%\end{figure}
%}
%
%\newcommand{\FigOptimStopTgtPosSensitivityIntegralRatio}{
%\begin{figure}[tb]
%\centering 
%%	\fbox{
%\includegraphics[width=0.9\textwidth,trim=1.2cm 0.05cm 0.9cm 0.6cm,clip]{figs/optimisation/StopTgtPosition/Tidied_-AcceptedMomentum-Integrated_ratios.pdf}
%%}
%\caption{\figlabel{optim:StopTgtPos:AcceptedIntegratedRatio}
%The relative avvep
%}
%\end{figure}
%}

\newcommand{\FigOptimDIOBeamBlockGeometry}{
\begin{figure}[tb]
\centering 
%\fbox{
\includegraphics[width=0.8\textwidth]{figs/optimisation/BeamAndDIOBlocker/StopTgt_to_Detector-cropped.png}
%}
\caption{\figlabel{optim:DIOBeamBlock:Geometry}
Location of the beam blocker and one possible geometry for the \ac{DIO} blockers, both highlighted in green, shown here before optimisation.
}
\end{figure}
}

\newcommand{\FigOptimDIOBeamBlockESTDispersion}{
\begin{figure}[tb]
\centering 
%\fbox{
\includegraphics[width=0.8\textwidth,trim=0.3cm 0 1.7cm 0.6cm,clip]{figs/optimisation/BeamAndDIOBlocker/Tidied_ElectronDispersion-EST.pdf}
%}\\\fbox{
\includegraphics[width=0.8\textwidth,trim=0.3cm 0 1.7cm 0.6cm,clip]{figs/optimisation/BeamAndDIOBlocker/HandTidied_ElectronDispersion-EST-envelope.pdf}
%\includegraphics[width=0.8\textwidth,trim=1.5cm 0 8.5cm 3.0cm,clip]{figs/optimisation/BeamAndDIOBlocker/Tidied_ElectronDispersion-EST-envelope.png}
%}
\caption{\figlabel{optim:DIOBeamBlock:ESTDispersion}
Momentum-dependent dispersion of electrons passing through the spectrometer.
Top plot: the mean height of different momenta electrons as a function beamline distance, showing how the drift of the centre of gyration is truly proportional to the momentum.
Bottom plot: single standard deviation bands for electrons at different momentum, which shows how the envelope for different momenta overlap considerably, reducing the effectiveness of any collimators.
%Whilst the drift of the centre of gyration is proportional to momentum (top plot), there is reasonable overlap in the overall enevelope
%Mean height of electrons with different momentum from the stopping target to the detector before adding any DIO blockers and before optimised of the beam blocker.
}
\end{figure}
}

\newcommand{\FigOptimDIOBeamBlockAcceptances}{
\begin{figure}[tb]
\centering 
%\fbox{
\subfloat[\figlabel{optim:DIOBeamBlock:Acceptance:40}40 MeV/c]  {\includegraphics[width=0.35\textwidth,trim=0.0cm 0.3cm 0.5cm 1.5cm,clip]{figs/optimisation/BeamAndDIOBlocker/Tidied_Acceptance2D_40.pdf}}\hspace{1cm}
\subfloat[\figlabel{optim:DIOBeamBlock:Acceptance:60}60 MeV/c]  {\includegraphics[width=0.35\textwidth,trim=0.0cm 0.3cm 0.5cm 1.5cm,clip]{figs/optimisation/BeamAndDIOBlocker/Tidied_Acceptance2D_60.pdf}}\\
\subfloat[\figlabel{optim:DIOBeamBlock:Acceptance:80}80 MeV/c]  {\includegraphics[width=0.35\textwidth,trim=0.0cm 0.3cm 0.5cm 1.5cm,clip]{figs/optimisation/BeamAndDIOBlocker/Tidied_Acceptance2D_80.pdf}}\hspace{1cm}
\subfloat[\figlabel{optim:DIOBeamBlock:Acceptance:100}100 MeV/c]{\includegraphics[width=0.35\textwidth,trim=0.0cm 0.3cm 0.5cm 1.5cm,clip]{figs/optimisation/BeamAndDIOBlocker/Tidied_Acceptance2D_100.pdf}}
%}
\caption{\figlabel{optim:DIOBeamBlock:Acceptance}
Acceptance into the straw tracker for electrons with different momentum at the stopping target as a function of the beam and DIO blocker dimensions.
Note the logarthmic scale for the colour bar.
}
\end{figure}
}

\newcommand{\FigOptimDIOBeamBlockHitRate}{
\begin{figure}[tb]
\centering 
%\fbox{
%\hspace{-1.3cm}
\subfloat[\figlabel{optim:DIOBeamBlock:HitRate}Hit Rate from DIO]               {\includegraphics[width=0.45\textwidth,trim=0.8cm 0.3cm 0.3cm 1.05cm,clip]{figs/optimisation/BeamAndDIOBlocker/HitRate2d.pdf}}\hspace{0.04\textwidth}
\subfloat[\figlabel{optim:DIOBeamBlock:HitVAccept}Signal Acceptance vs Hit Rate]{\includegraphics[width=0.45\textwidth,trim=0.8cm 0.3cm 0.3cm 1.05cm,clip]{figs/optimisation/BeamAndDIOBlocker/HitRateVsAcceptance2dRatio.pdf}\vspace{1cm}{}}
%\subfloat[\figlabel{optim:DIOBeamBlock:HitVAccept}Signal Acceptance vs Hit Rate]{\includegraphics[width=0.45\textwidth,trim=0.8cm 0.3cm 2.5cm 1.05cm,clip]{figs/optimisation/BeamAndDIOBlocker/HitRateVsAcceptance2d.pdf}\vspace{1cm}{}}
%\hspace{0.12\textwidth}
%}
\caption{\figlabel{optim:DIOBeamBlock:HitRateAcceptance}
\protect\subref{fig:optim:DIOBeamBlock:HitRate} Number of straw tracker hits per DIO electron.
\protect\subref{fig:optim:DIOBeamBlock:HitVAccept} Ratio between the high-momentum electron ($p>100$~MeV/c) acceptance to the number of hits per DIO electron.  Colour is on a logarithmic scale.
}
\end{figure}
}

\chapter{Phase-II Optimisation}
\sectlabel{phaseII-optimisation}
%\begin{easylist}
%# Before a substantial sensitivity estimate can be made, need a solidly optimised design
%# Aiming for \sensePII within a single year of running
%# Designs previously optimised \cite{CDR}, and these results are used as nominal design / starting point
%# Fresh optimisation using new software / simulation, updated fieldmaps, physics lists and geometry
%# Some aspects fixed already since \phaseI under construction: Experiment hall, Torus1, detector solenoid, fieldmap and coil parameters?
%# Key areas for optimising:
%## Production target dimensions
%## Torus1 dipole field strength
%## Torus2 dipole field strength
%## Electron spectrometer dipole field strength
%## collimator shapes and locations
%## stopping target and beam blocker position and form
%## DIO blockers on spectrometer
%\end{easylist}
The last study into the sensitivity of \COMET \phaseII was performed in 2009~\cite{CDRphase2}, before the staged approach and \phaseI design had even been considered.
That study found that a \ac{ses} of $2.6\times10^{-17}$ could be achieved in $2\times10^{7}$ seconds of running with a total expected background count rate of fewer than 0.34 events over the entire run period.
Since then, the focus of the collaboration has shifted to the \phaseI design and R\&D and the \phaseII design has not been optimisised or studied further.

The purpose of this chapter is to apply the updates in the fieldmap calculation, the geometry handling and physics modelling to revisit the design of \phaseII and demonstrate that it can do at least as well as the previous 2009 design.
Some aspects of the design have been fixed already by \phaseI such as the experiment hall and the coils and cold-mass of the first stages of the muon beamline.
The aspects of the experiment that remain open for optimisation are shown in \tab{optimisation:possible-parameters}.

\TabOptimisationParameters
As an initial configuration, much of the design from the 2009 study will be used, with updates included for the areas of the experiment that have since been refined by the \phaseI design.

\section{Optimisation Strategy}
%\begin{easylist}
%# Take some aspects as fixed
%# Limit scope and approach:
%## Ideally, each aspect optimised in combination to maximise signal acceptance and reduce background
%## How decoupled are each section?
%## In practise such an optimisation is not easy to do, instead aim to produce a baseline optimisation so that all backgrounds / issues can be identified
%## This can then form basis for further optimisation, with perhaps a smarter more integrated approach
%# Method:
%## Production target optimisation
%### Maximise muon and pion yield between 0 and 80 MeV at entrance to muon beamline
%### Parameters to vary: target length, target radius
%## Muon beam optimisation
%### Maximise muon stopping rate in stopping target
%### Minimise pion stopping rate
%### vary dipole along TS2 and TS4
%### vary Collimators: TS2 and at TS3
%## Electron spectrometer optimisation
%### Optimise dipole to increase signal acceptance
%### Optimse DIO blockers so DIO rate per straw is less than 1 kHz
%### Vary solenoidal field to increase separation?
%## Stopping target / beam blocker optimisation
%### Maximise reflection of signal electrons from upstream by tuning target position
%## Detector optimisation
%\end{easylist}

Some 32 parameters are listed in \tab{optimisation:possible-parameters} and yet this number alone does not describe the full challenge of optimising the \phaseII design since many of these parameters will be correlated.
It would seem natural for correlations to exist between the muon beam collimator(s), the stopping target and the beam blocker, since all involve removing or stopping muons and other particles in the beam.
Other less intuitive correlations might also exist however, and a complete optimisation would need to include the impact of these as well.

This optimisation study then can be considered a scan through a (greater than) 32 dimension parameter space.
A brute force search of such a space would be nightmarishly slow and require an enormous amount of computing power.
Machine learning algorithms or intelligent scanning techniques might be able to tackle such a problem, and perhaps in the future these methods will be used.
In the meantime however I am forced to use the technique of 'physicists intuition' to approach the problem, whereby some parameters are assumed uncorrelated whilst others are disregarded on the expectation that their impact be small.

The goal of this chapter therefore is not only to optimise the experiment but also to evaluate the correlations.
ideally this will also identify which parameters are the most strongly correlated, which would help make future optimisation studies more efficient.
The outputs of this optimisation should not be considered as final but should instead be treated as a baseline from which more intelligent approaches or physicists can improve.

%\section{Optimisation Goals}
%\begin{easylist}
%# Set sensitivity goal and optimise to reach this
%# Identify which parameters are strongly correlated to help focus and improve the efficiency of future optimisation studies.
%# Single event sensitivity only considers signal acceptance, but also need to understand backgrounds in terms of final confidence limit that can be set
%\end{easylist}

\section{Production Target Optimisation}
In the \phaseII \ac{CDR}, the production target is given as being 16~cm in
length and 4~mm in radius~\cite{CDR}.  Since then, there have been changes to
the magnetic field in this region, as well as the lengths and locations of
solenoids, shielding and beam-pipe, and the proton beam.  Previous studies have
looked at comparing the Tungsten target proposed for Phase-II to other
materials~\cite{thesis-AEdmonds}, and also drawn a comparison between
MARS~\cite{MARS}, Geant4~\cite{Geant4} and the limited data available.

The goal in this study then is to optimise the production target with the
up-to-date configurations.  This study aims to maximise the total muon and pion
yield below 80~MeV at the entrance to the Torus1 bent solenoid, by varying the
radius and length of the production target.

%\begin{easylist}
%# Some studies in the past (Andy's thesis) have compared Tungsten with Graphite and optimised Phase-I target
%# Geometry in CDR was 16~cm length and 4~mm radius
%# Proton beam, capture solenoids and radiation shielding has changed since then, so re-optimisation of the production target was needed
%\end{easylist}

\subsection{Configuration} 
%\icedustDescription{one}{tow}
%begin{easylist}
% Configuration
%## Only using Geant4 in ICEDUST
%# ICEDUST description: 
%## \icedustDescription{heads/1512w51_develop(3a0ee59)__3_UNCOMMITTED__}{heads/Patch_Geant4-G4MultiLevelLocator(11fc8f0)}
%# Production target:
%## Material: Tungsten
%## Back of target fixed 8~cm behind nominal centre (intersection of Proton beam axis and muon beam-line axis)
% Proton beam configuration
%## Uncertainty over beam profile makes this exercise difficult
%# Treat incoming beam as a pencil beam which is not realistic, but reasonable for this study \CHECK{Is it? Need to quantify this}
%# Use proton beam parameters from TDR2014:
%## Double Gaussian with horizontal spread of 5.8 mm and vertical spread of 2.9 mm.
%## Energy distribution also Gaussian with mean kinetic energy 8.01 GeV and 0.135~MeV spread in sigma.
%## The beam was fed in 1~cm upstream of target front face, which is the largest cause for errors in this optimisation study.
%end{easylist}

\Tab{optimisation:ProdTgtSec:configuration} gives the key parameters for the
beam input and other aspects of this simulation.  The location and orientation
of the target were held fixed, since the proton beamline is fixed with respect
to the muon beam axis.  Once a realistic proton beam becomes available, these
values would also benefit from optimisation, however.  During the scan over
length, the back face of the target was kept 8~cm away from the muon beam since
the radiation shielding has previously been optimised, and since beyond this
the magnetic field will no longer be able to capture the pions and muons
produced.

It must be noted that at this point in time there is an appreciable
uncertainty in the proton beam profile and position.  In particular, whilst the
proton beamline upstream has been delivered, the effect of the magnetic
field and necessary dipole and quadrupole magnetics are still being studied by
the proton beam-line group.  The beam profile is given in the \phaseI \ac{TDR}
as having a Gaussian profile and energy distribution, but no divergence or
location is given.  The effect of the proton beam distribution on the overall
sensitivity shall therefore be considered later on.

Protons originated from a plane (distributed as a two dimensional Gaussian
across this surface) but since there is therefore some scope to tune the proton
beam's position, the input particle plane was moved to remain 1~cm away from
the front surface of the target.  Since the aim is to maximise the muon and
pion yield by varying only the length and radius, shifting the proton beam
input plane in this way removes any variation of target acceptance due to
divergences of the proton beam in the magnetic field.

\begin{table}[t]
\centering
\begin{tabular}{lp{0.7\textwidth}}
	\hline
	\multicolumn{2}{l}{\bf Proton Beam}\\
	Horizontal spread, $\sigma_x$ & 5.8 mm \\
	Vertical spread, $\sigma_y$ & 2.9 mm \\
	Mean energy, $\mu_E$ & 8.01 GeV \\
	Energy spread, $\sigma_E$ & 0.135 MeV \\
	\hline
	\multicolumn{2}{l}{\bf Target}\\
	Material & Tungsten  \\
	Orientation & 10\degree between target's principal axis and the muon beam axis.    \\
	Location & Back face fixed 8~cm away from muon beam axis.  \\
	Length & 16~cm in CDR. Varied in steps of 4~cm from 4 to 32~cm.\\
	Radius & 4~mm in CDR. Varied from 2 to 10~cm in steps of 2~cm and from 10 to 30~cm in steps of 4~cm.  \\
	\hline
	\multicolumn{2}{l}{\bf Software configuration}\\
	Packages & \ttfamily{heads/1512w51_develop(3a0ee59)__3_UNCOMMITTED__}  \\
	Externals & \ttfamily{heads/Patch_Geant4-G4MultiLevelLocator(11fc8f0)}  \\
	Fieldmap & 160104 \CHECK{} \\
	\hline
	\multicolumn{2}{l}{\bf Sample Sizes}\\
	Length scan &  3e5 \ac{POT} (30 runs of 1e4) \\
	Radius scan & 4.9e5 \ac{POT} (49 runs of 1e4)  \\
	Final scan &   \\
\hline
\end{tabular}
\caption{Key parameters in the configuration of the Production Target optimisation.}
\tablabel{optimisation:ProdTgtSec:configuration}
\end{table}

\subsection{Length Scan}
Different length targets were simulated with 3e5 POT per length.
Target length was varied in steps of 4~cm from 4 to 32~cm, whilst the target radius was held fixed at the CDR value of 4~mm.

\Fig{optimisation:ProdTgtSec:Length:Momentum} shows the momentum distributions of pions and muons for different target lengths.
Target length is given as half-length which is the Geant4 convention.  
\Fig{optimisation:ProdTgtSec:Length:Integral} then shows these distributions integrated up to different momentum.
From these plots it can be seen that for both muons and pions, the optimum target length occurs around a total length of 32~cm.

Additionally it can be seen from
\fig{optimisation:ProdTgtSec:Length:IntegralRatio} that the shape of the
momentum distributions changes only weakly as a function of the target length.
These plots were produced by normalising the integrated momentum contours of
\fig{optimisation:ProdTgtSec:Length:Integral} to the total integral below
400~MeV.  As a result, it is possible that the actual shape variation is even
weaker than apparent here, since in the present sample size, the high momentum
tail is not well sampled at small target lengths, such that a skew in the
normalisation might occur.

\FigOptimProdTgtLength

\subsection{Radius scan}
In parallel to the length optimisation scan, different radii targets were also simulated.
Targets with radii of 2, 4, 6, 8, 10, 14, 18, 22, 26, and 30~mm were tested.
The target length was held at the \ac{CDR} value of 16~cm in total.

The results of these scans can be seen in \fig{optimisation:ProdTgtSec:Radius:Momentum} and \fig{optimisation:ProdTgtSec:Radius:Integral},
where it can be seen that a maximum in both the muon and pion yields at the entrance to the Torus1 section is achieved at a radius of about 10~mm.
As in the length scan, the shape variation of the momentum distributions is rather weak as a function of target radius.

\FigOptimProdTgtRad

\subsection{Final Result}
Since the length and radius scan were performed in parallel, a final cross check was performed where the optimal radius was confirmed at the optimised target length.
The integrated spectrum is shown in \fig{optimisation:ProdTgtSec:Combined:Integral} where it can be seen that the optimum radius once the target length is increased to 32~cm is still 10~mm.

\begin{easylist}
	# Figure comparing Phase-II to Phase-I
	# Conclusions
\end{easylist}

\begin{figure}[t]
\centering
\includegraphics[width=0.45\textwidth,trim=0 0 0 1.5cm,clip]{figs/optimisation/ProdTgtGeom/Radius_pi-minus_integral_ratios}
\caption{
Comparison of the combined muon and pion yield at the entrance to Torus1 for \phaseI (blue) and \phaseII (red).
}
\figlabel{optimisation:ProdTgtSec:ComparePhI-II}
\end{figure}

\section{Dipole Strengths of the Muon Beamline}
%\begin{easylist}
%# Production target simulation
%	## Plots of input distribtion (possibly shown in previous section)
%	## Mention gas pressure at exit of ProdTgtSec but show why this shouldn't have been an issue
%# Optimisation technique
%	## Vary dipole field strength in Torus1 and Torus2 simultaneously
%	## Mesh scan through parameter space
%	## Look at muon stopping rate
%	## Stopping target as in the CDR 
%	## No collimator material included
%# Results
%	## Muon stopping rate vs dipole field plot
%	## Pion stopping rate vs dipole field plot
%	## Discussion of corraltion
%# Selected optimal value
%\end{easylist}

The full 180\degree bent solenoid that makes up the bulk of the muon transport beam line is actually broken down into two
90\degree pieces -- Torus1 and Torus2 (also known as TS2 and TS4 by the magnet group).
Each of these sections has its own dipole field which need not provide the same dipole field strength.
The Torus1 section has been built already and was previously optimised for both \phaseI~\cite{TDR2016} and \phaseII~\cite{CDRphase2}. 
The dipole coils that it contains are designed to produce a dipole field of 0.055~T.
By running a lower current through these coils one could reduce the Torus1 dipole field without too much effort. 
However if \phaseII should require a greater dipole field strength for this region
that might be trickier since it could require additional windings to be inserted -- this would not be impossible but might costly.

\subsection{Large-sample Production Target Simulation}
Before the muon beam section could be optimised, a large set of \ac{POT} events were produced transporting all particles up to the entrance of the muon beam section.
In total, $2.4\times10^{8}$~POT were simulated through this stage, equivalent to about 1.4 bunches at \phaseII.
\CHECK{Are these numbers correct?}
The production target used the optimised geometry from the previous section.
All particles that hit the surface of the Torus1 container volume were read out for later re-use.
In addition, particles entering the proton beam dump were also saved if later simulations wished to study their impacts.

\CHECK{plot that shows the momentum distributions of particles in this input}

\begin{table}[t]
\CHECK{Fill this table in properly}
\centering
\begin{tabular}{lp{0.7\textwidth}}
	\hline
	\multicolumn{2}{l}{\bf Pions and muons from Prod.\ Tgt.\ simulation}\\
	%Energy spread, $\sigma_E$ & 0.135 MeV \\
	\hline
	\multicolumn{2}{l}{\bf Software configuration}\\
	Packages & \ttfamily{---}  \\
	Externals & \ttfamily{---}  \\
	Fieldmap & 160104 \CHECK{} with additional scale factors applied to Torus dipole field\\
	\hline
	\multicolumn{2}{l}{\bf Sample Sizes}\\
	Final scan &   \\
\hline
\end{tabular}
\caption{Key parameters in the configuration of the muon beamline dipole field optimisation.}
\tablabel{optimisation:ProdTgtSec:configuration}
\end{table}

\subsection{The Optimised Dipole Field Strengths}
The figure-of-merit for this optimisation is the muon stopping rate, which will be maximised at the optimal field configuration.
%Toshiba have previously performed a realistic field calculation using OPERA for a single coil whose on-axis field has a mean of about 0.055~T.
To identify such a configuration, a 2D-grid scan was done where the Torus1 and Torus2 dipole fields were simultaneously varied.
Varying the field strength was done by applying a scale factor to the fieldmaps for each 90\degree section described in \sect{sw:fieldmap}.
These scale factors were varied in steps of 0.125 (equivalent to 6.875~mT) and for each point muons and pions from $8\times10^6$~POT were transported to the stopping target.
The stopping rate for each combination of scale factors was then assessed.

\FigOptimMuBeamDipoleMuStops
\FigOptimMuBeamDipolePiStops

The muon stopping rate as a function of the two dipole field strengths are shon in \fig{optim:muBeamDipole:stoppedMu}.
A scale factor of 1 on the x-axis means no change to the current \phaseI design and since there is little improvement by moving to larger values for Torus1 the optimal scale factors are chosen to be 1 and 0.5 for the Torus1 and TTorus2 respectively.  
This translates to dipole field strengths of 0.055 and 0.0275~T respectively.
\CHECK{Are these the final selected dipole field strengths?}

Also striking from \fig{optim:muBeamDipole:stoppedMu} is the anti-correlation between the two dipole field strengths.  
Roughly speaking the sum of the optimal dipole field values is constant, \ie $B_1^\text{optimal} + B_2^\text{optimal} = \text{const}$.
Although this was not previously foreseen, such a correlation can be understood accordingly.
The stopping target geometry and position is held fixed in this parameter scan, which essentially fixes the upper momentum of muons that can stop in the target.
%The initial momentum distribution grows roughly linearly with momentum for the momenta of muons that will stop in the target (see the later sections for more on this), which would suggest tuning to the higher momentum would increase the number of muons reaching the target.
%That effect however is greatly reduced by the fact that the initial muon beam has a very large transverse position profile
The drift of a particle due to the dipole field is proportional only to the distance travelled in that dipole field and the dipole field strength, and does not depend on the particle's momentum.
This means that the total drift due to the Torus1 and Torus2 dipole components is proportional to the weighted sum of the two dipole field strengths, where the weight is the distance travelled through each of the two sections.
Since each section is the same length, the time for a particle to pass through Torus1 will be roughly the same as the time in Torus2 so that the weights would be roughly equal.
This leads to the total drift being proportional to the sum of the two dipole fields.
Since the production target (as the source of muons) and stopping target (as the muons' destination) are at a fixed height with respect to one another the optimal vertical drift is also fixed so that the sum of the dipole field strengths should also be roughly fixed.

With that said, the Torus1 section has a higher pion flux which causes some assymetry between the two sections.
Keeping more pions on-axis in the Torus1 section means that more muons will enter the Torus2 section from those pions that have decayed.
But since the pion momentum distribution is slightly higher than the muon distribution keeping pions on axis requires a larger dipole field strength.
This could explain the slight assymetry where the muon stopping rate appears slightly larger if Torus1's field is larger than that of Torus2.

It is also interesting to consider the pion stopping rate as a function of the dipole field strengths.
However, as can been from \fig{optim:muBeamDipole:stoppedPi} the stopping rate is close to the level of POT events used in the simulation so that the plot is dominated by statistical fluctuations.

\section{Collimators in the Muon Beamline}
\begin{easylist}
# Muon beam without collimators
	## Height vs momentum before and after 180 degrees
	## Beamline projection for stopped and high-P muons
	## Subtraction of the good and bad muons to suggest the region of interest for collimators
	## Why the assymetry upstream?
# Collimator optimisation
	## Analytic collimator study to reduce simulation time
	### Cannot study secondary particles
	### Assumes all particles entering the collimators are killed
	## Current collimator from Phase-I are in a different location to those suggested here
# Results of study, recommended collimator design
\end{easylist}

\section{Electron Spectrometer's Dipole}

\section{Stopping Target Position}
\section{The Beam and Decay-in-Orbit Blockers}

\section{Revisiting the Torus2 Collimator after Beam Background Studies}
\section{Summary of optimised parameters}
\section{Future optimisations}
\begin{easylist}
# Stopping target shape, thickness, disk spacing etc
# High energy electron acceptance vs beamline distance
# Correlation between dipole fields and stopping target shape
# Simultaneous signal efficiency and background acceptance optimisation (use final expected confidence limit as figure of merit)
\end{easylist}
